\section{Example 3: Kinematics}

For the final example, we give a more complete simulation from a classical
kinematics.  This example is selected both because it is familiar to the reader and because it demonstrates what happens when the systems we model increase in complexity.  

Consider shooting a cannon ball off of a cliff. We ask where will the ball land below?
We start by modeling this mathematical problem in SymPy without statistics:

\begin{figure}[ht]
\vspace{-0pt}
\centering
\includegraphics[width=.55\textwidth]{images/kinematics}
\vspace{-0pt}
\caption{Variables involved in shooting a cannon ball from a cliff}
\label{fig:kinematics}
\vspace{00pt}
\end{figure}

\begin{lstlisting}
    # Symbols for position, velocity, angle, time, and gravity
    >>> x0, y0, yf, v, theta, t, g = symbols('x0 y0 yf v theta t g')
    # x and y positions as a function of time
    >>> x = x0 + v*cos(theta)*t
    >>> y = y0 + v*sin(theta)*t + g/2*t**2
    # Solve y = yf for time to obtain the duration of the flight
    >>> t_impact = solve(y-yf, t)[1]
    # Final x value of cannon ball on impact
    >>> xf = x0 + v*cos(theta)*t_impact
\end{lstlisting}
This code yields the following symbolic solution:

\begin{equation}
\label{eqn:kinematics_symbolic}
x_f = x_{0} - \frac{v \left(v \sin{\left (\theta \right )} + \sqrt{- 2 g y_{0} + 2 g y_f + v^{2} \sin^{2}{\left (\theta \right )}}\right) \cos{\left (\theta \right )}}{g}
\end{equation}

To obtain numeric results for a particular set of values, we can replace the symbolic variables, $x_0, y_0, y_f, v, \theta, g$ with numeric ones:

\begin{lstlisting}
# Substitute out symbolic variables in xf for numeric values
>>> xf_num = xf.subs({ g: -9.8, v: 10, theta: pi/4, 
                       x0: 0, y0: 10, yf: 0})
\end{lstlisting}

\begin{equation}
\label{eqn:kinematics_numeric}
x_f = \frac{250}{49} + \frac{50}{49} \sqrt{123} = 16.4189
\end{equation}

The expression in Eqn. \ref{eqn:kinematics_symbolic} is less appealing to the
naked eye since it is hard to visualize its meaning but it serves two
purposes. Because it is a full symbolic description of the problem, it can
provide insight as to the relation of the variables to the
result. Additionally, it serves as a staging ground to produce many numeric
solutions such as the one in Eqn. \ref{eqn:kinematics_numeric}. Once the
mathematical model has been described no further conceptual work is necessary
on the part of the programmer to produce numeric results.

\subsection{Adding Uncertainty}

We now consider the statistical problem. Imagine first that we are uncertain
about the initial height of the cannon. Rather than knowing the precise value
$y_0 = 10$, we suspect it is distributed normally with standard deviation 1, 
$y_0 \sim \N{10}{1}$. Due to this uncertainty $x_f$ no longer has a precise
value. Instead of asking for an exact number, we may ask for the
expectation value, variance, or even the full distribution.

Traditionally, this is done either through painstaking algebra (for simple
cases) or by building external codes to sample using Monte Carlo methods. We
now show how SymPy can be used to model this statistical problem and provide
solutions.

Rather than defining $y_0$ as a symbolic or numeric variable, we declare
it to be a random variable distributed normally with mean 10 and standard
deviation 1.
\begin{lstlisting}
>>> y0 = Normal(10,1)
\end{lstlisting}

This variable acts in all ways like a SymPy symbol and so we may still form the
symbolic expression in Eqn. \ref{eqn:kinematics_symbolic}. 
Because this expression now contains random elements, it itself is random and is active under our statistical operators.

\begin{lstlisting}
>>> E(xf)
\end{lstlisting}

This code follows standard theory to generate the following integral
$$\int_{-\infty}^{\infty} \frac{25}{49} \frac{\left(\sqrt{\frac{98}{5} x_{1} + 50} + 5 \sqrt{2}\right) e^{- \frac{1}{2} \left(x_{1} -10\right)^{2}}}{\sqrt{\pi}}\, dx_{1}$$

This integral is challenging and the solution includes special
functions. Fortunately, algorithmic tools are capable of producing exact
analytic solutions. The full solution is omitted here in the interests of space and clarity; numeric equivalents are presented below in their place (floating point results are requested using the {\tt evalf()} method). 

$$\textrm{E}(x_f)\textrm{.evalf()} = 16.4099...$$
$$\textrm{var}(x_f)\textrm{.evalf()} = .2044...$$

In principle, we could ask for the full analytically computed density of $x_f$
although the complexity prohibits any meaningful insight.

\subsection{Dealing with Complexity}

Imagine now a more difficult problem where more of the variables in our model were uncertain
\begin{lstlisting}
>>> v = Uniform(9,11.5) # Uniformly distributed between 9 and 11.5
>>> theta = Normal(pi/4, pi/20)
\end{lstlisting}

Our queries may also become more complex such as, \textit{What is the
probability that theta was greater than $\pi/4$ given that the ball landed
past $x_f>16$?} , $P(\theta>\frac{\pi}{4} | x_f>16)$. Asking this question in
SymPy is just as easy as the example above:
\begin{lstlisting}
>>> P(theta>pi/4  ,  xf>16)
\end{lstlisting}
Solving the resultant integral however is substantially more challenging. For
situations like these exact solutions are often infeasible and Monte Carlo
codes are built around the problem. Because this problem statement is now
cleanly separated from the solution code, it is trivial to connect a simple
sampling engine onto SymPy stats allowing all such queries to be approximated
using a simple syntax
\begin{lstlisting}
>>> P(theta>pi/4,  xf>16, numsamples=10000)
0.3428
\end{lstlisting}

The path to solution presented here is typical in a lab environment. First a mathematical model is built, then initial numeric results are computed and finally errors, uncertainty and statistics are addressed. The work presented here shows how a computer algebra system that includes symbolic statistics can facilitate this process in a real-world setting.
